\section{Introduction}\label{sec:p2-introduction}
\subsection{From the MEA benchmark to transfer across amine classes}
% P:p2-introduction-01
A shared electrolyte description is useful only to the extent that it describes observations beyond those used to estimate it.
The companion Part I develops an explicit-ion electrolyte perturbed-chain statistical associating fluid theory (ePC-SAFT) description for loaded aqueous monoethanolamine (MEA), with equilibrium, caloric and density properties calculated from a common thermodynamic potential \cite{PolleyPartI}.
Its pressure AARDs at unfitted MEA concentrations are 33.89\,\% at 15 wt\% and 44.53\,\% at 45 wt\%, despite closer agreement at the fitted concentration.
Those errors motivate a test of the shared description across other amines; they do not supply a common accuracy standard for different experimental data.

% P:p2-introduction-02
Here the water, \COtwo and inorganic-ion parameters from Part I remain fixed while only amine-specific quantities are fitted.
MDEA is a tertiary amine whose reaction network contains bicarbonate formation without amine carbamate.
AMP is a sterically hindered primary amine, DEA a secondary amine, and PZ a cyclic diamine with mono- and dicarbamate chemistry.
These choices test the common bicarbonate description both without carbamate and alongside different carbamate pathways.
One example per class cannot isolate chemical class from differences in observation coverage, species representation or fitting freedom.

\subsection{Calibration, unfitted observations and phase behavior}
% P:p2-introduction-03
Published SAFT amine models distinguish parameter estimation from predictions to different extents.
Baygi and Pahlavanzadeh fit pure and binary MEA properties and then calculate ternary solubility using ideal-activity chemical speciation \cite[Sects.~3.1--3.3]{Baygi2015}.
Nasrifar and Tafazzol also organize MEA, DEA and MDEA around pure properties, binary interactions and subsequent acid-gas comparisons, using an ideal-activity treatment rather than explicit electrolyte activity coefficients \cite[Sects.~2--6]{Nasrifar2010}.
Wangler et al. estimate MDEA ion--water interactions using osmotic data before predicting reactive ternary equilibria \cite[Sects.~6.1--6.2.1]{Wangler2018}.
The present fits instead include amine-specific reactive observations, so agreement with those observations is calibration.

% P:p2-introduction-04
In this paper, held-out observations never enter a fitting objective.
That term does not imply that the data remained unseen throughout selection of the final fit.
Previously accessed observations outside the fit and comparisons made after selecting the parameters are identified separately.
Finite searches for additional phases provide numerical evidence about the calculated liquids, while experimental pressure and property comparisons measure agreement with the named observations.
Neither form of evidence alone establishes independent validation of ion parameters or qualification for absorber design.
The results follow the properties being tested, so fitted pressures, unfitted pressures, other thermodynamic properties and phase behavior each constrain the same transfer hypothesis.
